\chapter{Esquema de la base de datos SQLite}
\label{anexo:bbdd}

La base de datos SQLite empleada para la persistencia de datos de diagnóstico se
compone de cuatro tablas relacionadas. El esquema completo, incluyendo los índices
y la configuración WAL, se muestra a continuación.

\section*{D.1 — Esquema SQL}

\begin{lstlisting}[language=SQL, caption={Esquema completo de la base de datos diagnostics.db}]
PRAGMA journal_mode=WAL;

-- Sesiones de diagnostico
CREATE TABLE IF NOT EXISTS sessions (
    id         INTEGER PRIMARY KEY AUTOINCREMENT,
    label      TEXT    NOT NULL DEFAULT '',
    started_at TEXT    NOT NULL,  -- ISO 8601 UTC
    ended_at   TEXT               -- NULL si la sesion sigue activa
);

-- Muestras de telemetria OBD-II
CREATE TABLE IF NOT EXISTS samples (
    id           INTEGER PRIMARY KEY AUTOINCREMENT,
    session_id   INTEGER NOT NULL REFERENCES sessions(id),
    pid          INTEGER NOT NULL,   -- PID OBD-II en decimal
    name         TEXT    NOT NULL,   -- nombre legible del parametro
    value        REAL    NOT NULL,   -- valor decodificado con unidades
    unit         TEXT    NOT NULL,   -- unidad de medida (rpm, km/h, ...)
    monotonic_ts REAL    NOT NULL,   -- time.monotonic() para ordenacion
    wall_ts      TEXT    NOT NULL    -- ISO 8601 UTC para presentacion
);

-- Comandos ejecutados y tramas CAN asociadas
CREATE TABLE IF NOT EXISTS commands (
    id           INTEGER PRIMARY KEY AUTOINCREMENT,
    session_id   INTEGER NOT NULL REFERENCES sessions(id),
    command      TEXT    NOT NULL,   -- nombre del comando (snapshot, dtcs, vin...)
    request_hex  TEXT    NOT NULL,   -- trama ISO-TP de peticion en hex
    response_hex TEXT    NOT NULL,   -- trama ISO-TP de respuesta en hex
    wall_ts      TEXT    NOT NULL    -- ISO 8601 UTC
);

-- Codigos de averia (DTC) leidos en cada sesion
CREATE TABLE IF NOT EXISTS dtcs (
    id           INTEGER PRIMARY KEY AUTOINCREMENT,
    session_id   INTEGER NOT NULL REFERENCES sessions(id),
    code         TEXT    NOT NULL,   -- codigo J2012 (p.ej. P0420)
    raw_hex      TEXT    NOT NULL,   -- 2 bytes crudos del DTC en hex
    description  TEXT    NOT NULL DEFAULT '',
    wall_ts      TEXT    NOT NULL    -- ISO 8601 UTC
);

-- Indices para consultas por sesion
CREATE INDEX IF NOT EXISTS idx_samples_session
    ON samples(session_id, pid);

CREATE INDEX IF NOT EXISTS idx_commands_session
    ON commands(session_id);

CREATE INDEX IF NOT EXISTS idx_dtcs_session
    ON dtcs(session_id);
\end{lstlisting}

\section*{D.2 — Descripción de las tablas}

La tabla \textbf{sessions} registra cada sesión de diagnóstico con sus timestamps de
inicio y fin. El campo \texttt{ended\_at} es \texttt{NULL} mientras la sesión está
activa y se actualiza al cerrarla mediante \texttt{end\_session()}.

La tabla \textbf{samples} almacena cada muestra individual de telemetría. El campo
\texttt{pid} permite filtrar las muestras de un PID concreto dentro de una sesión,
y el campo \texttt{request\_hex} / \texttt{response\_hex} en la tabla
\texttt{commands} almacena las tramas CAN en formato hexadecimal compacto (sin
espacios), lo que permite reconstruir la secuencia ISO-TP completa para cada
comando ejecutado.

La tabla \textbf{commands} registra todos los comandos ejecutados desde la aplicación
móvil, incluyendo los de gestión de DTCs y obtención de VIN. El campo
\texttt{response\_hex} de un comando de modo 0x09 (VIN) contiene la concatenación
de la First Frame y las Consecutive Frames recibidas.

La tabla \textbf{dtcs} almacena los códigos de avería leídos en cada sesión, con el
código en formato J2012 (\texttt{code}), los dos bytes crudos del DTC (\texttt{raw\_hex})
y, cuando está disponible, su descripción textual. Se rellena al ejecutar una lectura
de DTCs (modo 0x03).

\section*{D.3 — Ejemplo de contenido}

La tabla~\ref{tab:ejemplo_samples} muestra un fragmento representativo del contenido
de la tabla \texttt{samples} tras una sesión de monitorización de cuatro PIDs.

\begin{table}[!htbp]
\caption{Fragmento de la tabla \texttt{samples} — sesión de ejemplo.}
\label{tab:ejemplo_samples}
\centering
\renewcommand{\arraystretch}{1.2}
\begin{tabular}{|c|c|l|c|l|l|}
\hline
\textbf{id} & \textbf{pid} & \textbf{name} & \textbf{value} & \textbf{unit} & \textbf{wall\_ts} \\
\hline
1 & 12  & rpm            & 1724.0 & rpm  & 2025-04-10T10:23:01Z \\
2 & 13  & vehicle\_speed & 87.0   & km/h & 2025-04-10T10:23:01Z \\
3 & 5   & coolant\_temp  & 91.0   & °C   & 2025-04-10T10:23:02Z \\
4 & 17  & throttle\_pos  & 23.5   & \%   & 2025-04-10T10:23:02Z \\
5 & 12  & rpm            & 1756.0 & rpm  & 2025-04-10T10:23:02Z \\
\hline
\end{tabular}
\end{table}

\section*{D.4 — Consideraciones de concurrencia}

La activación del modo WAL (\textit{Write-Ahead Logging}) mediante
\texttt{PRAGMA journal\_mode=WAL} es esencial para el correcto funcionamiento
concurrente del sistema. Sin WAL, las escrituras del hilo de monitorización
bloquearían las lecturas del hilo BLE durante el volcado del buffer, introduciendo
latencias variables en las respuestas a la aplicación móvil. Con WAL, los lectores
no se bloquean por los escritores en ningún caso, al operar sobre versiones
consistentes del fichero de base de datos.

El buffer de escritura de 50 muestras implementado en \texttt{SqliteDataLogger}
reduce la frecuencia de escrituras a disco a una transacción cada 50 muestras
(aproximadamente cada 25\,segundos a 500\,ms de periodo de muestreo con cuatro
PIDs activos), minimizando el tiempo de bloqueo del \textit{write lock} WAL y
el desgaste de la tarjeta SD de la Raspberry Pi.
